% Editorial supplement; not one of the 187 translated source components.
\clearpage
\renewcommand{\CurrentExpose}{译校注}
\phantomsection
\addcontentsline{toc}{section}{译校注}
\section*{译校注}

以下说明由译校者增补, 不属于原刊正文. 原刊的陈述、证明和原编号保留;
正文的“校记”标志指向本附录. 本轮独立 AI 对读覆盖指定英文读者版的全部
187 个组件, 但不等于专业人工审定, 也不等于法文原刊逐行对勘已经完成.

\medskip
\noindent\hypertarget{editorial-vi-mixedchar}{}\underline{校记 A：VI 6.7 的底环特征与判别式.}

VI 第 6 节只规定 $\Lambda$ 为剩余域为 $k$ 的完备诺特局部环.
6.7 又假设 $\operatorname{char}k=2$, 却没有要求 $2=0$ 于 $\Lambda$.
因此不能在一般底环上略去所写形变方程对 $X_0$ 求导时的 $2X_0$ 项.
原刊全书页 124、131--132 (档案 PDF 第 129、136--137 页) 确有上述假设和平方公式.

对 6.7 证明给出的显式模型, 令
\[
 R=\Lambda[[t,t_1]],\qquad
 F=X_0^2+\sum_{i=1}^{n}X_{2i-1}X_{2i}+tX_0+t_1,
 \qquad A=R[[X_0,\ldots,X_{2n}]]/(F).
\]
相对 Jacobian 理想实际为
\[
 I_A(A|R)=(2X_0+t,X_1,\ldots,X_{2n}).
\]
所以, 记 $X=X_0$, 相应有限 $R$-模为
\[
 C=A/I_A(A|R)
   \simeq R[[X]]/(X^2+tX+t_1,\,2X+t).
\]
环 $B=R[[X]]/(X^2+tX+t_1)$ 是以 $1,X$ 为基的秩二自由 $R$-模.
这里二次式首一, 其低次系数属于完备局部底环的极大理想, 可用 Weierstrass 除法.
在此基下, 乘以 $2X+t$ 的矩阵为
\[
 M=\begin{pmatrix}t&-2t_1\\2&-t\end{pmatrix},
 \qquad C\simeq\operatorname{Coker}(M:R^2\longrightarrow R^2),
\]
因为 $(2X+t)X=-2t_1-tX$. 按 VI 5.9 所用的零次 Fitting 理想定义, 得
\[
 I_R(A|R)=\operatorname{Fitt}^0_R(C)
         =(\det M)=(t^2-4t_1).
\]

若 $2=0$ 于 $\Lambda$, 则这确实恢复原刊的 $(t^2)$.
原刊所写 $C\simeq R/(t)\oplus R/(t)$ 也在这个条件下成立.
但取 $\Lambda=\mathbb Z_2$, 原刊的一般平方断言不成立;
这并不是选择另一组形式参数就能排除的问题.
事实上, 若在 $\mathbb Z_2[[t,t_1]]$ 中有
\[
 t^2-4t_1=u h^2,\qquad u\in R^*,
\]
将其映入 $\mathbb Q_2[[t,t_1]]$ 后, 左端最低总次数为 $1$,
右端则为 $2\operatorname{ord}(h)$, 是偶数, 矛盾.
两个非零主理想若相等, 其生成元在此整环中相差单位,
故上述论证也排除了把该理想写成某个形式参数平方所生成的理想.

另一方面, 原模型中可消去 $t_1$, 仍有
\[
 A\simeq\Lambda[[t,X_0,\ldots,X_{2n}]],
\]
所以这个模型在 $\Lambda$ 上的形式光滑性不受上述订正影响.
应区分形式光滑性和 Jacobian 理想为参数平方这两个结论.
6.8(2) 所引的局部平方公式同样须加上底环特征为 $2$ 的限制,
不能仅凭剩余域特征为 $2$ 及原证明将其推广到混合特征.
此处给出了显式模型的订正计算, 没有另行断言 6.8 的所有一般化均已证明.

\medskip
\noindent\hypertarget{editorial-vi-square-root}{}\underline{校记 B：VI 6.3 的单位平方根步骤.}

6.3(1) 的证明由 $k=k^2$ 及 $R$ 完备推断每个所需单位均为平方.
这一步在特征 $2$ 时不能只由所列条件得到: 对特征 $2$ 的完美域 $k$,
$R=k[[s]]$ 中的单位 $1+s$ 不是平方, 因为平方幂级数的奇数次系数均为零.
若 $2$ 可逆, 可将剩余域中的平方根用 Hensel 引理提升;
保留原刊特征 $2$ 的适用范围时, 则仍需另外证明所需二次型变换的存在性.
这一例子定位了证明步骤的缺口, 不构成对 6.3 同构判据本身的反例.
本译本没有以该反例为由删去命题或擅补新的假设.

其余已识别的原刊假设、符号与引文问题, 以及本轮明确采用的记名订正,
逐项列于源码包的审校记录. 这些记录将原刊疑点与中文误译分开,
不把忠实于有疑点的底本称作数学论证已经全部审定.
